% GENERATED FILE. Edit canonical lessons, then run npm run generate:books.
% canonical-source-hash: fnv1a64:446be2233a890dae
% canonical-lessons: MR-C409-khalasi, MR-C409-ohoti, MR-C409-udhan, MR-C409-khavalne, MR-C409-hindkalne

\chapter{Sailors and Tides}
\label{ch:mr-sailors-and-tides}
\hlchaptermodality{409}

\begin{chapteropening}
\textbf{By the end of this chapter:} \emph{I can talk about sailors, low tide, a swelling sea and a boat rocking.}

Complete the last lesson of chapter 409: I can talk about sailors, low tide, a swelling sea and a boat rocking.
\end{chapteropening}

\section[\texorpdfstring{\mr{खलाशी} (khalāśī)}{खलाशी (khalāśī)}]{\mr{खलाशी} (khalāśī) — a sailor}

\label{lesson:MR-C409-khalasi}



\begin{quote}
Before the new one: say the Marathi for to toss about, then the Marathi for an anchor.
\end{quote}

\subsection*{You'll want to know: \mr{खलाशी}}
\textbf{\mr{खलाशी}} — \emph{khalāśī} — \textquotedblleft{}a sailor\textquotedblright{}.

\begin{quote}
\mr{जहाजावर} \mr{वीस} \mr{खलाशी} \mr{काम} \mr{करतात}.
\end{quote}

\begin{quote}
\emph{jahājāvar vīs khalāśī kām kartāt.}
\end{quote}

\begin{quote}
Twenty sailors work on the ship.
\end{quote}

Masculine.

The plural is the same in the direct case: \mr{एक} \mr{खलाशी}, \mr{वीस} \mr{खलाशी}.

\begin{cousinweb}
\textbf{\mr{खलाशी}} (\emph{khalāśī}) is a loan from Persian, usually traced to \emph{khalāṣī}, a word for a ship's hand.
\end{cousinweb}

\begin{grammarlens}[title={What you've built}]
One more word for getting around.
\end{grammarlens}

\subsection*{Your turn}
\begin{itemize}
\raggedright
  \item \emph{Say it:} \emph{khalāśī}
  \item \emph{Say it:} \emph{khalāśī}, once more
  \item \emph{Say it:} \emph{jahājāvar vīs khalāśī kām kartāt.} — Twenty sailors work on the ship.
  \item \emph{Recall:} say the Marathi for to toss about, then the Marathi for an anchor, then say \emph{khalāśī} again
\end{itemize}

\begin{tcolorbox}[breakable,colback=teal!4,colframe=teal!35!black,title={Before you move on}]
What does \mr{खलाशी} mean? (\textquotedblleft{}A sailor\textquotedblright{}.) Say it once more.
\end{tcolorbox}
\section[\texorpdfstring{\mr{ओहोटी} (ohoṭī)}{ओहोटी (ohoṭī)}]{\mr{ओहोटी} (ohoṭī) — low tide; an ebb}

\label{lesson:MR-C409-ohoti}



\begin{quote}
Before the new one: say the Marathi for an anchor, then the Marathi for a sailor.
\end{quote}

\subsection*{You'll want to know: \mr{ओहोटी}}
\textbf{\mr{ओहोटी}} — \emph{ohoṭī} — \textquotedblleft{}low tide; an ebb\textquotedblright{}.

\begin{quote}
\mr{ओहोटीला} \mr{वाळूवर} \mr{खेकडे} \mr{दिसतात}.
\end{quote}

\begin{quote}
\emph{ohoṭīlā vāḷūvar khekaḍe distāt.}
\end{quote}

\begin{quote}
At low tide you can see crabs on the sand.
\end{quote}

Feminine.

Time words take -\mr{ला}: \mr{ओहोटीला}, at low tide.

\begin{cousinweb}
\textbf{\mr{ओहोटी}} (\emph{ohoṭī}) is inherited, from the family of \emph{*avaghaṭṭati}, \textquotedblleft{}decreases, recedes\textquotedblright{}, where CDIAL lists Marathi \emph{ohaṭṇẽ}, \textquotedblleft{}to abate, ebb\textquotedblright{} (CDIAL 748).
\end{cousinweb}

\begin{grammarlens}[title={What you've built}]
One more word for getting around.
\end{grammarlens}

\subsection*{Your turn}
\begin{itemize}
\raggedright
  \item \emph{Say it:} \emph{ohoṭī}
  \item \emph{Say it:} \emph{ohoṭī}, once more
  \item \emph{Say it:} \emph{ohoṭīlā vāḷūvar khekaḍe distāt.} — At low tide you can see crabs on the sand.
  \item \emph{Recall:} say the Marathi for an anchor, then the Marathi for a sailor, then say \emph{ohoṭī} again
\end{itemize}

\begin{tcolorbox}[breakable,colback=teal!4,colframe=teal!35!black,title={Before you move on}]
What does \mr{ओहोटी} mean? (\textquotedblleft{}Low tide; an ebb\textquotedblright{}.) Say it once more.
\end{tcolorbox}
\section[\texorpdfstring{\mr{उधाण} (udhāṇ)}{उधाण (udhāṇ)}]{\mr{उधाण} (udhāṇ) — a sea surge; a spring tide}

\label{lesson:MR-C409-udhan}



\begin{quote}
Before the new one: say the Marathi for a sailor, then the Marathi for low tide.
\end{quote}

\subsection*{You'll want to know: \mr{उधाण}}
\textbf{\mr{उधाण}} — \emph{udhāṇ} — \textquotedblleft{}a sea surge; a spring tide\textquotedblright{}.

\begin{quote}
\mr{आज} \mr{समुद्राला} \mr{मोठं} \mr{उधाण} \mr{आलं}.
\end{quote}

\begin{quote}
\emph{āj samudrālā moṭhaṁ udhāṇ ālaṁ.}
\end{quote}

\begin{quote}
Today the sea surged high.
\end{quote}

Neuter.

\mr{उधाण} \mr{येणे}, a surge comes; it also names a rush of joy: \mr{आनंदाला} \mr{उधाण} \mr{आलं}.

\begin{cousinweb}
\textbf{\mr{उधाण}} (\emph{udhāṇ}) is inherited: CDIAL lists Marathi \emph{udhāṇ}, \textquotedblleft{}spring-tide\textquotedblright{}, under \emph{uddhāna}, \textquotedblleft{}rising up\textquotedblright{} (CDIAL 2015).
\end{cousinweb}

\begin{grammarlens}[title={What you've built}]
One more word for getting around.
\end{grammarlens}

\subsection*{Your turn}
\begin{itemize}
\raggedright
  \item \emph{Say it:} \emph{udhāṇ}
  \item \emph{Say it:} \emph{udhāṇ}, once more
  \item \emph{Say it:} \emph{āj samudrālā moṭhaṁ udhāṇ ālaṁ.} — Today the sea surged high.
  \item \emph{Recall:} say the Marathi for a sailor, then the Marathi for low tide, then say \emph{udhāṇ} again
\end{itemize}

\begin{tcolorbox}[breakable,colback=teal!4,colframe=teal!35!black,title={Before you move on}]
What does \mr{उधाण} mean? (\textquotedblleft{}A sea surge; a spring tide\textquotedblright{}.) Say it once more.
\end{tcolorbox}
\section[\texorpdfstring{\mr{खवळणे} (khavaḷṇe)}{खवळणे (khavaḷṇe)}]{\mr{खवळणे} (khavaḷṇe) — to grow rough}

\label{lesson:MR-C409-khavalne}



\begin{quote}
Before the new one: say the Marathi for low tide, then the Marathi for a sea surge.
\end{quote}

\subsection*{You'll want to know: \mr{खवळणे}}
\textbf{\mr{खवळणे}} — \emph{khavaḷṇe} — \textquotedblleft{}to grow rough\textquotedblright{}.

\begin{quote}
\mr{संध्याकाळी} \mr{समुद्र} \mr{अचानक} \mr{खवळला}.
\end{quote}

\begin{quote}
\emph{sandhyākāḷī samudra acānak khavaḷlā.}
\end{quote}

\begin{quote}
In the evening the sea suddenly grew rough.
\end{quote}

Past (he …): \mr{तो} \mr{खवळला}, \emph{to khavaḷlā}.

For the sea and for people: \mr{तो} \mr{खवळला}, he flared up.

\begin{cousinweb}
\textbf{\mr{खवळणे}} (\emph{khavaḷṇe}) is a Marathi verb for the sea, or a temper, becoming wild; its origin is not settled.
\end{cousinweb}

\begin{grammarlens}[title={What you've built}]
One more word for getting around.
\end{grammarlens}

\subsection*{Your turn}
\begin{itemize}
\raggedright
  \item \emph{Say it:} \emph{khavaḷṇe}
  \item \emph{Say it:} \emph{khavaḷṇe}, once more
  \item \emph{Say it:} \emph{sandhyākāḷī samudra acānak khavaḷlā.} — In the evening the sea suddenly grew rough.
  \item \emph{Recall:} say the Marathi for low tide, then the Marathi for a sea surge, then say \emph{khavaḷṇe} again
\end{itemize}

\begin{tcolorbox}[breakable,colback=teal!4,colframe=teal!35!black,title={Before you move on}]
What does \mr{खवळणे} mean? (\textquotedblleft{}To grow rough\textquotedblright{}.) Say it once more.
\end{tcolorbox}
\section[\texorpdfstring{\mr{हिंदकळणे} (hindkaḷṇe)}{हिंदकळणे (hindkaḷṇe)}]{\mr{हिंदकळणे} (hindkaḷṇe) — to slosh about; to rock}

\label{lesson:MR-C409-hindkalne}



\begin{quote}
Before the new one: say the Marathi for a sea surge, then the Marathi for to grow rough.
\end{quote}

\subsection*{You'll want to know: \mr{हिंदकळणे}}
\textbf{\mr{हिंदकळणे}} — \emph{hindkaḷṇe} — \textquotedblleft{}to slosh about; to rock\textquotedblright{}.

\begin{quote}
\mr{लाटांवर} \mr{होडी} \mr{हिंदकळत} \mr{होती}.
\end{quote}

\begin{quote}
\emph{lāṭāṁvar hoḍī hindkaḷat hotī.}
\end{quote}

\begin{quote}
The boat was rocking on the waves.
\end{quote}

Past (he …): \mr{तो} \mr{हिंदकळला}, \emph{to hindkaḷlā}.

For water moving in a pot or a tank, and for a boat on waves.

\begin{cousinweb}
\textbf{\mr{हिंदकळणे}} (\emph{hindkaḷṇe}) may belong with Sanskrit \emph{hindolayati}, \textquotedblleft{}swings\textquotedblright{}, where CDIAL lists Marathi \emph{hĩḍaḷṇẽ}, \textquotedblleft{}to jog, shake\textquotedblright{} (CDIAL 14095); the link is not settled.
\end{cousinweb}

\begin{grammarlens}[title={What you've built}]
One more word for getting around.
\end{grammarlens}

\subsection*{Your turn}
\begin{itemize}
\raggedright
  \item \emph{Say it:} \emph{hindkaḷṇe}
  \item \emph{Say it:} \emph{hindkaḷṇe}, once more
  \item \emph{Say it:} \emph{lāṭāṁvar hoḍī hindkaḷat hotī.} — The boat was rocking on the waves.
  \item \emph{Recall:} say the Marathi for a sea surge, then the Marathi for to grow rough, then say \emph{hindkaḷṇe} again
\end{itemize}

\begin{tcolorbox}[breakable,colback=teal!4,colframe=teal!35!black,title={Before you move on}]
What does \mr{हिंदकळणे} mean? (\textquotedblleft{}To slosh about; to rock\textquotedblright{}.) Say it once more.
\end{tcolorbox}
